% Pursuing Stacks (cotes 134-2 to 134-8): archive cote 134-5, site page 91 (PDF page 92), typescript page 278.
% typescript. Draft compared: ps4.tex:4666-4799.
% AI-transcribed by Claude Opus 5.5 (claude-opus-5-5, adaptive thinking, effort high), Message Batches API, batch ps2 (msgbatch_01NVwr1Yuo8qJf5rykeGBXPK), 2026-11-12; not yet checked by eye.
\scanpage{134-5/91}
\mspage{\hand{278}}

where B is a small category, \inserted{and moreover $i$ is a functor $i: A\to B\subset B^\wedge$.} This then brings us to the related

\uline{Question} 1': Let i: $A\to B$ be an aspherical map in (Cat), under which
restrictive conditions on A,B,i (if any) does \xout{\unsure{$i_!$}\illegible{1}}
\[ i_! : A^\wedge\to B^\wedge \]
take aspheric objects into aspheric ones, i.e. do we have
\[ (7) \qquad i_!(A^\wedge_{as}) \hand{$\subset$} B^\wedge_{as} \qquad ? \]

The question makes a sense even without assuming i to be aspheric.
If we drop this asphericity assumption on i, we are reduced, for a given
B , (\xout{and}\inserted[typed]{by} replacing A by $A_{/F}$ ) where F is a given \xout{object} aspheric object of
$A^\wedge$) to the case when A is aspheric and $F=e_{A^\wedge}$ , i.e. to looking at whe-%
ther the element
\[ (8) \qquad i_!(e_{A^\wedge}) \;=\; \varinjlim_{A}{}_{(B^\wedge)}\, i(a) \]
in $B^\wedge$\hand{/}is aspheric. This element appears as the direct limit in $B^\wedge$ of
aspheric elements i(a), the indexing category A being itself aspheric.
Thus it is rather tempting to hope that the limit might well be aspheric
too. Let's assume $\underline{W}=\underline{W}_\infty$ = usual weak equivalence, it turns out that when
A is 1-connected so is (8), which would seem to give some support to the
hope that (8) is aspheric if A is. However this is defintely not so, un-%
less at the very least we assume B to be equally aspheric (which is howe-%
ver a rather natural assumption, as it follows from A aspheric if we assu%
- me moreover i \inserted{to be} aspheric). To see this, we take i to be cofibering with
0-connected fibers - the cofibration assumption implies that the direct
limit\deleted{s} \inserted{over A} can be computed first fiberwise and then take a limit over B
(the most general version of associativity for direct limits ! ), whereas
the 0-connectedness assumption on the fibers implies that the limit taken
over the fiber $A_b$ is just b itself, hence (8) is isomorphic to
\[ \varinjlim_{b}\unsure{$(B^\wedge)$}\,{}^{b} \;=\; e_{B^\wedge} \quad , \]
which is'nt aspheric except precisely when B is; now for topological
reasons it is easy to find an aspheric A cofibered with 0-connected fibers
over a non-aspheric B.

Thus in the question of asphericity of (8) we'll better assume both
A \uline{and} B aspheric. If either A has a final object or B has an initial ob-%
ject, asphericity of (8) is more or less trivial in any case. For cer-%
tain categories B even without initial object, (8) is always aspheric
when A is, this I checked at any rate for the ordered category
\[ B \;=\; \diagram{\begin{tikzcd}[row sep=small, column sep=small] \hand{$\alpha$} \arrow[dr, "" ] & \\ & \hand{$\gamma$} \\ \hand{$\beta$} \arrow[ur] & \end{tikzcd}} \quad , \]

\margin{91}
